% Chapter 0. The rules. Shared by every building.

\section*{Chapter 0. The rules}
\phantomsection\addcontentsline{toc}{section}{\texorpdfstring{Chapter 0. The rules}{Chapter 0. The rules}}
\begin{center}
{\small What may be built, what a blur is, and what each line of a chapter claims.}
\end{center}

This chapter is the only place where the method itself is stated. Every later chapter uses these rules and adds none.

\subsection*{Background}

\begin{lines}
\item[\lab{S1}] A \emph{bag} is a set. We use sets, membership, subsets and maps in the ordinary way, and ordinary logic.
\item[\lab{S2}] A \emph{floor} is a bag $X$ together with a list $O$ of \emph{operations} (maps $X \times \dots \times X \to X$) and a list $Q$ of \emph{questions} (maps from $X$, or from $X \times X$, to some bag of answers; a yes/no question is a question).
\item[\lab{S3}] A \emph{law} is a statement about a floor. A law is an \emph{equation} if it has the form ``$s = t$ for all values of the letters'', where $s$ and $t$ are built from the letters and the operations only.
\end{lines}

\subsection*{Build}

Six ways to make a new bag from bags already made. Nothing else may be built.
\begin{lines}
\item[\lab{M1}] (Pairs and lists.) From $X$ and $Y$: the bag $X \times Y$ of pairs. From $X$ and $n$ in $\N$: the bag $X^n$ of lists of length $n$.
\item[\lab{M2}] (Sequences.) From $X$: the bag of all maps $\N \to X$.
\item[\lab{M3}] (Words.) From a list of names, each with a number of inputs: the smallest bag that holds $f(w_1, \dots, w_k)$ whenever $f$ takes $k$ inputs and it holds $w_1, \dots, w_k$. Two words are equal only if written the same. A map from the words to a floor is fixed by saying what it does to each name (definition by recursion).
\item[\lab{M4}] (Select.) From $X$ and a condition: the bag of those $x$ in $X$ that meet the condition.
\item[\lab{M5}] (Bags of bags.) From $X$ and a blur $\sim$ of $X$ (L1): the bag $X/{\sim}$ whose members are the bags of $\sim$.
\item[\lab{M7}] (Sub-bags.) From $X$: the bag of all sub-bags of $X$. \\
A \emph{map} $A \to X$ is a sub-bag of $A \times X$ in which each $a$ in $A$ stands on the left of exactly one pair. So the bag of all maps $A \to X$ is built by M1, M7, M4; M2 is the case $A = \N$. A blur of $X$ is a sub-bag of $X \times X$, so the bag of all blurs of $X$ is built the same way. M7 is the one line whose result is far larger than what went in; a chapter that uses it must say so.
\end{lines}
\begin{lines}
\item[\lab{M6}] (Lifting.) An operation of $X$ may be lifted to a built bag \emph{place by place}; any other operation on a built bag must be given by a formula in a Definition line. A selected bag keeps an operation only if it is closed under it, and the chapter must say so.
\end{lines}

\subsection*{Blur}

\begin{lines}
\item[\lab{L1}] (Blur.) A \emph{blur} of $X$ is a relation $\sim$ on $X$ with $x \sim x$; \ $x \sim y \Rightarrow y \sim x$; \ $x \sim y,\ y \sim z \Rightarrow x \sim z$. The \emph{bag} of $x$ is $\bag{x} := \{y : y \sim x\}$.
\item[\lab{L2}] (Respects.) $\sim$ \emph{respects} an operation $f$ if $x_1 \sim y_1, \dots, x_k \sim y_k \Rightarrow f(x_1, \dots, x_k) \sim f(y_1, \dots, y_k)$.
\item[\lab{L3}] (Keeps.) $\sim$ \emph{keeps} a question $q$ if $x \sim y \Rightarrow q(x) = q(y)$. For a question of two things: $x \sim x',\ y \sim y' \Rightarrow q(x,y) = q(x',y')$.
\item[\lab{L4}] (Blurrier.) $\sim$ is \emph{blurrier} than $\sim'$ if $x \sim' y \Rightarrow x \sim y$.
\item[\lab{L5}] (From below.) Given a list $G$ of \emph{gluings} (pairs to be glued) and a list $O$ of operations: the blur \emph{forced by $G$ respecting $O$} is the least blurry blur that glues every pair in $G$ and respects every operation in $O$.
\item[\lab{L6}] (From above.) Given a list $Q$ of questions and a list $O$ of operations: the blur \emph{kept by $Q$ respecting $O$} is the blurriest blur that keeps every question in $Q$ and respects every operation in $O$.
\item[\lab{L7}] (A blur must say what it respects.) Every Blur line names its $O$. The same gluing with a different $O$ is a different blur.
\item[\lab{L8}] (One-way blur.) A \emph{one-way blur} of $X$ is a relation $\preceq$ with $x \preceq x$ and $x \preceq y,\ y \preceq z \Rightarrow x \preceq z$. It respects $f$ as in L2. Its \emph{two-way part}, ``$x \preceq y$ and $y \preceq x$'', is a blur, and $\preceq$ is an order on the bags of that blur. A blur is a one-way blur that also runs backwards.
\item[\lab{L9}] (Likeness.) A \emph{likeness} on $X$ is a relation $N$ with $x \mathrel{N} x$ and $x \mathrel{N} y \Rightarrow y \mathrel{N} x$. It need not chain. A blur is a likeness that chains.
\item[\lab{L10}] (Limits.) $O$ holds only operations with finitely many inputs (S2). A blur may also be asked to \emph{respect limits} for a stated nearness (Chapter 5): if $x_n \sim y_n$ for every $n$, $x_n$ comes near $x$ and $y_n$ comes near $y$, then $x \sim y$. The Blur line then has a slot \emph{Limits:} naming the nearness. A limit is never put in $O$, and a question is never put in $O$.
\end{lines}

\subsection*{Theorems about blurs}

Stated once here, used in every chapter.
\begin{lines}
\item[\lab{Z1}] (Two ends.) Equality is the least blurry blur of $X$. ``Everything in one bag'' is the blurriest. Both respect every operation.
\item[\lab{Z2}] (From below exists.) The blur of L5 exists. It is reached by a process: start from $G$; add $x \sim x$; add $y \sim x$ for each $x \sim y$; add $f(\dots x \dots) \sim f(\dots y \dots)$ for each $x \sim y$ and each $f$ in $O$; add $x \sim z$ for each chain $x \sim y \sim z$; repeat until nothing new appears.
\item[\lab{Z3}] (From above exists.) The blur of L6 exists. For questions of one thing, it glues $x$ and $y$ exactly when $q(c(x)) = q(c(y))$ for every $q$ in $Q$ and every \emph{context} $c$: a map $X \to X$ made from the operations in $O$ and fixed members of $X$ (for example $c(x) = x + z$).
\item[\lab{Z4}] (Sandwich.) If every gluing in $G$ is kept by every question in $Q$ in every context, then the blur from below is no blurrier than the blur from above.
\item[\lab{Z5}] (Recipe for Appears.) In Z4 the two blurs are \emph{equal} as soon as: for every $x \not\sim y$ (from below) there is a context $c$ and a $q$ in $Q$ with $q(c(x)) \neq q(c(y))$.
\item[\lab{Z6}] (Descent.) If $\sim$ respects $f$, then $f(\bag{x_1}, \dots, \bag{x_k}) := \bag{f(x_1, \dots, x_k)}$ is an operation on $X/{\sim}$. If $\sim$ keeps $q$, then $q(\bag{x}) := q(x)$ is a question on $X/{\sim}$.
\item[\lab{Z7}] (Equations are free.) An equation that holds on $X$ also holds: on $X^n$ and on sequences with operations lifted place by place; on any selected bag closed under the operations; and on $X/{\sim}$ for any blur respecting the operations.
\item[\lab{Z8}] (Everything else must be checked.) A law that is not an equation (one that says ``or'', ``$\neq$'', ``there is'', or uses a question) may be lost by Build, Select or Blur. Each chapter must prove it again or record its loss.
\item[\lab{Z9}] (Normal forms.) A \emph{normal form} for $\sim$ is a map $\nu : X \to X$ with $\nu(x) \sim x$, and $x \sim y \Rightarrow \nu(x) = \nu(y)$. Then $\nu(\nu(x)) = \nu(x)$, and each bag holds exactly one $x$ with $\nu(x) = x$.
\item[\lab{Z10}] (A likeness is a blur seen through a picture.) If $\sim$ is a blur of $X$ and $p : X \to Y$ is a picture, then ``$y \mathrel{N} y'$ iff $y = p(x)$, $y' = p(x')$ for some $x \sim x'$'' is a likeness on $Y$, and it may fail to chain. Every likeness arises this way. So likeness is not a new primitive: it is made from blur and picture.
\end{lines}

\subsection*{Arrows}

\begin{lines}
\item[\lab{W1}] An \emph{arrow} is a map $f : X \to Y$ between floors, with a list of the operations and questions it respects: $f(x + y) = f(x) + f(y)$, and so on.
\item[\lab{W2}] An arrow is an \emph{inside} if it is one-to-one, and a \emph{picture} if it reaches every member of $Y$.
\item[\lab{W3}] (Every blur is a picture.) $x \mapsto \bag{x}$ is a picture $X \to X/{\sim}$ respecting whatever $\sim$ respects.
\item[\lab{W4}] (Every arrow is a blur, then an inside.) For an arrow $f$, ``$x \sim y \iff f(x) = f(y)$'' is a blur of $X$ respecting what $f$ respects, and $f$ is the picture $X \to X/{\sim}$ followed by an inside $X/{\sim} \to Y$.
\end{lines}

\subsection*{Kinds of chapter}

Every chapter is one of seven kinds, or a few of them in a stated order. The second line under a chapter title names the kinds, in order, with these words and no others. A chapter claims a blur only if it has a Blur line. In every kind, a name comes only after the line that makes or pins its object (a Notation line), old names go to the Remarks, and every sign reused from an earlier floor passes C3 in an Arrow line.
\begin{lines}
\item[\lab{K1}] (\emph{blur}.) Makes one new floor $X/{\sim}$. Slots: the table below.
\item[\lab{K2}] (\emph{build}.) Makes a new bag by M1--M7 and blurs nothing. The Blur line reads ``none''; Survives and Dies become \emph{carried} and \emph{not carried} (C5).
\item[\lab{K3}] (\emph{select}.) Pins one object by what it must do. Slots: \emph{Tests} (the conditions, a Select M4); \emph{At most one}; \emph{At least one} (a Build, or ``closed'', said once); \emph{Read off} (laws proved from ``only one''). A select may ask a question of a normal form of an earlier blur and find it kept on some bags (C7).
\item[\lab{K4}] (\emph{arrow}.) Reads one arrow between floors already made. Slots: the arrow; what it respects and what dies through it (W1); its own blur (W4); inside or picture (W2).
\item[\lab{K5}] (\emph{count}.) Existence from a count, said openly: what is counted, the count, and why the count forces the object. Its lines are Theorems, not Survives lines.
\item[\lab{K6}] (\emph{table}.) Rows, tests, entries (F2--F4); a survey of all arrows or all blurs of one sort is a table. What is read off is a conjecture until proved (F6).
\item[\lab{K7}] (\emph{rerun}.) Runs an earlier chapter with a new Given. Slots: \emph{free} (equations, carried by Z7) and \emph{checked again} (every other law, Z8), each line citing the line it reruns.
\end{lines}

\subsection*{The lines of a blur chapter}

A blur chapter (K1) makes one new floor. Its lines, in order, and what each one claims:

\begin{center}
\renewcommand{\arraystretch}{1.35}
\begin{tabular}{@{}l l >{\raggedright\arraybackslash}p{8.6cm}@{}}
\textbf{Label} & \textbf{Line} & \textbf{Claim} \\ \hline
A or G & Axioms / Given & laws assumed (Chapter 1 only), or earned earlier and cited \\
D & Definitions & Build (M1--M5, M7), lifted operations (M6), questions. Adds no truth \\
B1 & Blur, from below & the gluings $G$ and the respected $O$ (L5) \\
B2 & Blur, from above & the questions $Q$ and the same $O$ (L6) \\
P & Bags & a test for $x \sim y$ that can be computed \\
P & Normal forms & a $\nu$ as in Z9 \\
P & Survives & $\sim$ respects an operation (L2) or keeps a question (L3) \\
P & Dies & $\sim$ does not keep a question: one witness $x \sim y$, $q(x) \neq q(y)$ \\
P & Appears & B1 $=$ B2, proved by Z4 and Z5, naming the separating context \\
N & Notation & names for $X/{\sim}$ and its members. Only after the blur \\
P & Arrow & an arrow old $\to$ new, inside or picture, and what it respects \\
T & Theorems & statements inside the new floor, in the ordinary form \\
R & Remarks & other blurs, other roads, links to other chapters \\
E & Exercises & \\
\end{tabular}
\end{center}

\begin{lines}
\item[\lab{C1}] (Earned Given.) Apart from the axioms of Chapter 1, every Given line cites the line of an earlier chapter that proves it.
\item[\lab{C2}] (Survives names only what is defined.) A Survives or Dies line may name only operations and questions defined on the floor being blurred.
\item[\lab{C3}] (One sign, four steps.) A sign such as $+$ may be reused on the new floor only after: \emph{lift} (D), \emph{survive} (P), \emph{descend} (Z6), and \emph{agree} (the Arrow respects it).
\item[\lab{C4}] (Dies is for traps.) Every Dies line follows from the Blur line. Write only those a reader might wrongly use.
\item[\lab{C5}] (Build-only chapters.) If the Blur line is ``none'', Survives and Dies are replaced by \emph{carried} and \emph{not carried}: which laws of the Given still hold on the built floor (Z7, Z8).
\item[\lab{C6}] (The first choice is open.) Which gluings or questions to start from is the author's choice. Everything after that choice is forced, and is proved.
\item[\lab{C7}] (Reserved words.) \emph{Survives} names a blur and an operation it respects (L2) or a question it keeps (L3), and nothing else. A question that takes one answer on every member of some bags, but not of all, is \emph{kept on those bags}; such a line is a Theorem. \emph{Glued} means ``in one bag of a blur'' and never ``near'' (Chapter 5, P5, T1: a blur that holds a rung of nearness has one bag). A theorem with a hypothesis is a Theorem, not a Survives line.
\item[\lab{C8}] (Imported.) A step that a later line uses and that the books do not prove is written as a Given line of the form \emph{(Imported: source.) statement. Used by: lines.} Each building lists its Imported lines in a ledger on its first page. A classical fact that no line uses stays in the Remarks, marked ``reported''.
\item[\lab{C9}] (Made from above.) When the only gluings one could list are the pairs that the question already glues, the Blur line says ``B1: none (made from above)'', and the chapter makes no Appears claim. Listing $e \sim 0$ for each $e$ in the bag of $0$, and letting the respected operations spread it (Z2), is a from-below line; listing every glued pair is not.
\end{lines}

\subsection*{The second floor}

\begin{lines}
\item[\lab{F1}] A chapter's Notation gives a \emph{system}; its Arrow gives an arrow between systems; its Theorems give yes/no facts about the system.
\item[\lab{F2}] A \emph{test} is a law that can be asked of every system. A \emph{table} has systems as rows, tests as columns, and the facts as entries.
\item[\lab{F3}] (Blur on a table.) For a set $U$ of tests: $K \sim_U L$ iff $K$ and $L$ answer every test in $U$ alike.
\item[\lab{F4}] (Class.) For a set $T$ of tests: $\operatorname{class}(T)$ is the set of rows passing every test in $T$. A blur asks what tests cannot tell apart; a class asks what tests pick out.
\item[\lab{F5}] (Closed.) A system is \emph{closed} for a chapter if running the chapter on it gives an Arrow that is both an inside and a picture.
\item[\lab{F6}] (The table suggests; rows refute.) A statement read off a table is a conjecture. It becomes a theorem by proof from the tests, or dies when a row is built that breaks it.
\item[\lab{F7}] (One way only.) Floor two reads floor one, never the reverse. A first-floor chapter cites only first-floor chapters; a class of the second floor only names what its rows already have. A Remark may point ahead, but no line may rely on what it points to.
\item[\lab{F8}] (Loop.) A chapter may build a system $X$ and blur it by the action of an earlier system $S$. A \emph{loop} is a theorem of that chapter saying that $S$ comes back: a question on the bags whose answers match $S$, respecting the operations named. Verdicts: \emph{closed} ($S$ comes back exactly), \emph{changed} (larger or smaller), \emph{wild} (no rung of nearness keeps the bags apart; then keep the moves, do not forget them). The return is a theorem of the later chapter: by F7 no earlier line may rely on it.
\end{lines}

\subsection*{Remarks}

\begin{lines}
\item[\lab{R1}] (What is assumed.) S1, M1--M5 and M7 are a small part of ordinary set theory. The book does not replace that foundation; it organises what is built on it.
\item[\lab{R2}] (Why two blurs.) From below is how a blur is \emph{made}: it tells you what is glued. From above is why it is \emph{right}: it tells you what the gluing is for. Appears is the claim that the two agree.
\item[\lab{R3}] (Why L7.) On $\R^n$ the choice of $O$ alone separates four different roads (Chapter 6).
\item[\lab{R4}] (Why Z8 matters.) ``No zero divisors'' is not an equation. It is lost by Blur in $\Z/6$ (Chapter 3) and by Build in $\R^2$ (Chapter 6).
\item[\lab{R5}] (Old names.) Blur respecting $O$: congruence. One-way blur: compatible preorder. Likeness: tolerance relation. From below: generated congruence. From above: syntactic congruence. Z7: the laws preserved by products, subalgebras and quotients. W4: the first isomorphism theorem. F4: a concept, in the sense of formal concept analysis.
\item[\lab{R6}] (What is not a rule.) A rule here may rely on nothing after Chapter 0 (F7 read backwards). So Amplify, a theorem about power sums whose proof uses series and poles (Chapters 17, 20, 33, 34), is not a rule: it lives in the field notes on the edge (Building 2's notes, N4, formerly Z11). The one-way growth blur it uses is Chapter 17's. F8 stays: it is a reading convention and relies on no chapter.
\end{lines}

\subsection*{Exercises}

\begin{lines}
\item[\lab{E1}] On $\N$ with the operation $x \mapsto x + 1$, find the blur forced by the single gluing $5 \sim 2$. How many bags are there?
\item[\lab{E2}] Prove Z4.
\item[\lab{E3}] Show that ``$x + y = y + x$'' is an equation and ``$x \neq 0 \Rightarrow$ there is $y$ with $x \times y = 1$'' is not. Find where the second is lost in this book.
\item[\lab{E4}] Give two different normal forms for the blur of Chapter 2.
\end{lines}
